% ============================================
% Supplementary Note 1: Outcome Classification, Metric Sensitivity, and Null Models
% ============================================

\subsection*{Supplementary Note 1: Outcome classification, metric sensitivity, and null models}
\label{sec:suppl:note1}

\rev{This Note formalizes outcome classification, tests metric sensitivity, and evaluates null models for passive abundance and additive effects.}

\paragraph{Outcome classification.}
\rev{To characterize coalescence outcomes, we used cosine similarity as the underlying compositional similarity metric between the post-coalescence community and each parental community:}
\begin{equation}
\rev{\text{Sim}(C, A) = \frac{\vec{x}_C \cdot \vec{x}_A}{\|\vec{x}_C\|_2 \|\vec{x}_A\|_2}, \quad
\text{Sim}(C, B) = \frac{\vec{x}_C \cdot \vec{x}_B}{\|\vec{x}_C\|_2 \|\vec{x}_B\|_2},}
\end{equation}
\rev{where $\vec{x}_A$, $\vec{x}_B$, and $\vec{x}_C$ denote community composition vectors, and $\|\cdot\|_2$ denotes the Euclidean norm. Because cosine similarity is invariant to multiplying either vector by a positive scalar, sample-wise total-abundance normalization does not change the similarity score. These cosine-similarity coordinates define the two-parent similarity space, where the position reflects how strongly the coalesced community resembles each parental community.}

Because parental communities may share species (overlapping ASVs), classification was based on a linear decomposition of the coalesced community in this normalized two-parent similarity space. We computed similarity-space projection coefficients $(u, v)$ from the weighted combination of parental compositions that reconstructs the coalesced community in this two-parent basis, and then retained only non-negative parental components because negative parental similarity has no biological interpretation. Thus, cosine similarity is the compositional metric, while vector decomposition provides the classification coordinates used to define retention and PDI. When parental communities share no species and the parental composition vectors are $L_2$ normalized, these coefficients reduce to the cosine similarities. The residual component not represented by either parental community quantifies novel restructuring in the same normalized space.

From these coefficients, we derived two classification metrics:

Retention magnitude ($r$): The coefficient-space retention score from both parental communities, calculated as $r = \sqrt{u^2 + v^2}$. Values close to 1 indicate high similarity-space retention of parental composition. Values close to 0 indicate extensive restructuring.

Parental Dominance Index (PDI): A measure of directional parental similarity, ranging from 0 (complete dominance by parental community B) through 0.5 (equal contributions) to 1 (complete dominance by parental community A). PDI is computed from the angular position in similarity space relative to the diagonal where both parental communities contribute equally:
\begin{equation}
\text{PDI} = \frac{2}{\pi} \operatorname{atan2}(u, v)
\label{eq:pdi}
\end{equation}
where $u$ and $v$ are the similarity-space projection coefficients for parental communities A and B, respectively. When $u \gg v$, PDI approaches 1 (parental community A dominance). When $v \gg u$, PDI approaches 0 (parental community B dominance). When $u = v > 0$, PDI equals 0.5 (equal contributions). Events with $u=v=0$ have zero retention and are classified as Restructuring rather than assigned a directional dominance interpretation.
\rev{Here $\operatorname{atan2}(u,v)$ denotes the two-argument arctangent using the common $\operatorname{atan2}(y,x)$ convention with $y=u$ and $x=v$. For non-negative $u$ and $v$, it returns an angle in $[0,\pi/2]$. When $v>0$ it is equivalent to the ratio-based one-argument arctangent, and when $v=0$ and $u>0$ it returns the boundary value $\pi/2$.}

Based on these metrics, each coalescence event was classified into one of three outcome categories. Restructuring occurs when $r^2 \leq 0.5$, \rev{equivalently when the retention radius satisfies $r \leq 1/\sqrt{2}$ in the two-parental-community similarity map,} meaning that parental-community retention is low in the normalized similarity space, indicating substantial ecological reorganization. Mixture occurs when $r^2 > 0.5$ and PDI is near 0.5 (i.e., $0.25 \leq \text{PDI} \leq 0.75$), indicating high retention with balanced similarity to both parental communities. Dominance occurs when $r^2 > 0.5$ and PDI is near 0 or 1 (i.e., $\text{PDI} < 0.25$ or $\text{PDI} > 0.75$), indicating high retention but with the composition of one parental community dominating that of the other.

We chose the classification thresholds based on two considerations. \rev{First, the retention threshold $r^2 = 0.5$ (radius $r = 1/\sqrt{2}$, not $r = 1/2$) defines a coefficient-space retention boundary in the two-parental-community similarity map. For the special case of non-overlapping, orthonormal parental composition vectors, this boundary corresponds to half of the squared similarity-space mass lying along the two parental axes. With overlapping parental vectors, it should be interpreted as an operational retention-score boundary rather than a literal explained-composition fraction.} Values below this indicate that novel species combinations dominate over parental similarity. Second, the PDI thresholds of 0.25 and 0.75 correspond to contribution ratios of approximately 3:1 between parental communities, a level of asymmetry that represents clear dominance of one parental community over the other while allowing for minor similarity to the subordinate parental community.

\rev{To make the boundary dependence explicit, we report the continuous similarity measures underlying the categorical classifications across media in Supplementary Figs.~10--12. A boundary-sensitivity analysis (Supplementary Fig.~37) further shows that the nutrient-dependent increase in Dominance remains evident when either the PDI boundary or retention boundary is varied around the manuscript baseline.}

\paragraph{Metric sensitivity.}
We compared coalescence outcome classifications using five approaches: the primary cosine-similarity/vector-decomposition framework, Euclidean distance, Bray-Curtis dissimilarity, Jensen-Shannon divergence, and Jaccard index. The alternative relative-abundance approaches, Euclidean distance and Bray-Curtis dissimilarity, produced outcome distributions broadly consistent with the primary framework, with Dominance as the most frequent outcome. However, Jensen-Shannon divergence and Jaccard index did not recover this pattern (Extended Data Fig.~2), reflecting the sensitivity of boundary assignments to the choice of compositional metric. We interpret this divergence as biologically informative rather than solely technical: the primary framework, Euclidean distance, and Bray-Curtis retain quantitative abundance information, whereas Jaccard collapses the data to presence/absence and therefore measures species-identity retention rather than abundance-weighted parental dominance.

\paragraph{Abundance-skew null models.}
A potential concern is that Dominance could arise from skewed abundance distributions rather than ecological interactions. To address this, we compared experimental parental-asymmetry values against two null models.

In Null Model 1 (abundance-weighted random selection), species survival probability is proportional to abundance in the combined parental-community pool, testing whether high-abundance species simply survive regardless of parental-community affiliation. In Null Model 2 (shuffled abundance), abundances are randomly permuted among species within each parental community before neutral mixing, testing whether the overall skewness structure alone drives Dominance.

We generated 500 null offspring communities per model and calculated parental asymmetry, $|2\mathrm{PDI}-1|$ (Supplementary Fig.~44). Experimental Base-medium data showed mean parental asymmetry of 0.698 ($n = 83$). Both null models produced significantly lower values: abundance-weighted null mean = 0.476 ($p < 10^{-18}$), shuffled null mean = 0.557 ($p < 10^{-11}$; Mann-Whitney U tests). These results indicate that abundance skewness alone cannot explain the observed patterns and support the conclusion that ecological processes beyond simple abundance-weighted survival help determine which parental community dominates.

\paragraph{\rev{Coalescence-pair-specific simple additive null model.}}
\rev{We also tested a stricter coalescence-pair-specific simple additive null model, $n_{C,\mathrm{null}} = n_A+n_B$, in which each null coalesced community was constructed from the raw ASV count vectors of the same two parental communities used in the corresponding experimental coalescence pair (Extended Data Fig.~3). Under this coalescence-pair-specific simple additive null model, an apparent Dominance classification can arise when the two parental-community raw ASV count vectors have substantially different Euclidean norms: for non-overlapping parental-community raw ASV count vectors, $\mathrm{Sim}(A,A+B)$ and $\mathrm{Sim}(B,A+B)$ are proportional to $\|n_A\|_2$ and $\|n_B\|_2$, respectively. This norm-imbalance effect is possible in principle, especially when richness is low and vector norms are sensitive to abundance unevenness among a small number of taxa. In the experimental data, however, the coalescence-pair-specific simple additive null model was usually classified as Mixture, including 66/83 Base medium events (79.5\%), and the observed Base medium outcomes were shifted further toward parental asymmetry than their corresponding null-model outcomes (one-sided paired Wilcoxon test, $p = 5.9 \times 10^{-14}$). Thus, the coalescence-pair-specific simple additive null model does not explain the observed Dominance tendency.}
